\documentclass[11pt]{article}
\usepackage[a4paper,margin=25mm]{geometry}
\usepackage{amsmath,amssymb,amsthm,hyperref}
\newtheorem{theorem}{Theorem}
\newtheorem{lemma}{Lemma}
\newcommand{\E}{\mathbb E}
\title{Polynomially many ordered rows with a fixed number of rational rank columns}
\author{Research proof: independently reviewed}
\date{30 September 2026}
\begin{document}\maketitle
\noindent\textbf{Status.} The full proof, including the continuous-density
baseline, has been independently reviewed by the root and rational-design
research agents.
\medskip

\noindent\textbf{Scope.} This is an auxiliary moment theorem for a fixed
number of rational cumulative-rank columns. It does not establish the
ordered-prefix identities or positive joint completion needed for three
or more simultaneous small dice.

\begin{lemma}[A lacunary fully fair continuous baseline]\label{fullbaseline}
Fix $u\ge0$ and $k\ge1$. Suppose positive integers $d_i$ satisfy
\[
 d_i>u+k-1+\sum_{j<i}d_j.
\]
For $|\eta_i|<1$, let the $i$th exceptional continuous variable have
density $f_i(t)=1-\eta_iP_{d_i}(2t-1)$ on $[0,1]$.
Together with $u$ independent uniform variables, these $k$ independent
variables are fully permutation-fair. All densities are strictly
positive; if the $\eta_i$ are rational, their coefficients are rational.
\end{lemma}
\begin{proof}
Fix one complete order and expand the product of the $k$ densities.
The empty correction term is its all-uniform probability $1/(u+k)!$.
For a nonempty correction subset, choose its largest index $i$ and
condition on that variable having value $t$. Integrating all other
ordered variables below and above $t$ gives an ordinary polynomial in
$t$ of degree at most
\[
 u+k-1+\sum_{j<i}d_j<d_i.
\]
Here each integration adds at most one to the total degree, and only
the selected lower-index perturbations contribute their degrees.
Legendre orthogonality makes the remaining integral against
$P_{d_i}(2t-1)$ zero. Thus every nonempty term vanishes. Positivity
and normalization follow from $|P_{d_i}|\le1$ and its zero mean.
\end{proof}

\begin{theorem}\label{main}
Fix integers $k\ge1$ and $T\ge1$. For every $m\ge1$, there exist
$K=O_k((m+1)^4)$ ordered real nodes $0<b_1<\cdots<b_K<1$ and $k$
strictly increasing rational columns $0<r_{i,1}<\cdots<r_{i,K}<1$ such that
\begin{equation}\label{mom}
 \frac1K\sum_{q=1}^K b_q^a\prod_{i\in S}r_{i,q}
       =\frac1{a+|S|+1},
 \qquad S\subseteq[k],\quad 0\le a\le m.
\end{equation}
All ranks have one common denominator at most
$C_{k,T}(m+1)^{C_{k,T}}$. The nodes have consecutive and endpoint gaps
at least $c_k/K$, and
\[
 \max_{i,q}|r_{i,q}-b_q|\le C_{k,T}(m+1)^{-T}.
\]
All constants are independent of $m$.
\end{theorem}

The proof uses two independently reviewed inputs: the separated
midpoint-grid quadrature lemma in
\path{../rational_designs/asymptotic/two_discrete_step_compatibility.tex}
and the polynomial conditioning theorem in
\path{../rational_designs/asymptotic/lacunary_legendre_conditioning.tex}.
For fixed $k,T$, below ``polynomial'' means bounded by a fixed power
of $m+1$, with constants and exponent allowed to depend on $k,T$.

\begin{proof}
Put $H_d(t)=\int_0^t P_d(2u-1)\,du$. It vanishes at both endpoints,
has derivative bounded by one, and is orthogonal in Lebesgue measure
to every polynomial of degree at most $d-2$. For example the latter
fact follows from its expression as a linear combination of shifted
Legendre polynomials of degrees $d-1$ and $d+1$.
Choose integers
\[
 d_i=m+k+2+\sum_{j<i}(d_j+1),\qquad 1\le i\le k,
\]
so all degrees are $O_k(m+1)$, and put
\[
 R_i(t)=t-\eta H_{d_i}(t),\qquad
 \eta=c(m+1)^{-(T+1)}
\]
for a fixed sufficiently small positive rational $c$.
These maps increase from zero to one with derivative at least $1/2$.
Expanding a product and integrating the factor with largest index in
each nonempty correction term gives
\begin{equation}\label{base}
 \int_0^1 t^a\prod_{i\in S}R_i(t)\,dt
 =\frac1{a+|S|+1},\qquad a\le m.
\end{equation}
Indeed the polynomial multiplying the largest $H_{d_i}$ has degree
at most $m+k+\sum_{j<i}(d_j+1)\le d_i-2$.

Let $D=m+\sum_i(d_i+1)=O_k(m+1)$. Choose the reviewed real
equal-weight quadrature exact through degree $2D$, with
$K=O_k((m+1)^4)$ nodes $b_q^0$ at distance at most $1/(8K)$ from
the uniform midpoint grid. Set $r^0_{i,q}=R_i(b_q^0)$.
All identities~\eqref{mom} hold at this real baseline, and all rank
and node gaps are at least $c_k/K$.
In fact this real baseline has a stronger full-order interpretation:
take the $k$ densities in Lemma~\ref{fullbaseline} with these degrees
and parameters, together with $m+1$ uniforms, and replace one uniform
by the quadrature variable $B$. Its conditional complete-order tests
have degree at most $m+k+\sum_i d_i=D$, so the quadrature preserves
all complete orders. The subsequent rational-rank construction below
retains the stated mixed threshold moments; no claim that it retains
this stronger full-order interpretation is made.

\medskip\noindent\emph{Rationally fixing the constant mixed moments.}
Choose an integer $L_0$ of sufficiently large polynomial size and
round all $r^0_{i,q}$ to multiples of $1/L_0$. This makes an error
at most $1/L_0$ in each entry. We next correct exactly the finitely
many identities~\eqref{mom} with $a=0$ and $S\ne\varnothing$.

Partition the row indices into $k$ interlaced groups $Q_1,\ldots,Q_k$.
Proceed with $j=k,k-1,\ldots,1$. At stage $j$, vary column $j$ only
on a fixed number $d_j^*=2^{k-j}$ of selected rows in $Q_j$, to impose
the equations for all $S$ whose least element is $j$. These equations
are exactly linear in the proposed changes; their coefficient rows
are the evaluations of
\[
 \prod_{i\in U}r_i,\qquad U\subseteq\{j+1,\ldots,k\}.
\]
Previously fixed equations do not involve column $j$ and are unchanged.
Because the row groups are disjoint, all other entries on the rows
of $Q_j$ still equal their initially rounded values.

We justify a polynomially conditioned square minor. At the baseline,
the corresponding continuous polynomial family is independent and
has polynomial least singular value by the reviewed lacunary rank-product
theorem. Sampling on $Q_j$ preserves this up to a constant: its star
discrepancy is $O_k(1/K)$, and for any polynomial $F$ of degree at most
$D$,
\[
 \left|\frac1{|Q_j|}\sum_{q\in Q_j}F(b_q^0)^2-\int F^2\right|
 \le C D^2\operatorname{disc}(Q_j)\|F\|_2^2.
\]
The right side is at most half the squared norm after the absolute
constant in $K$ is increased. Initial rounding preserves polynomial
coercivity if $L_0$ is a sufficiently large power. The Cauchy--Binet
formula now supplies a $d_j^*\times d_j^*$ evaluation minor with
polynomially bounded inverse: its dimension depends only on $k$,
and its entries are bounded. Solving that rational minor corrects
all equations at this stage with changes of size at most
$\operatorname{poly}(m)/L_0$.

The denominator increase is also polynomial. Each matrix entry has
denominator dividing a fixed power of $L_0$, its dimension is fixed,
and its determinant numerator is bounded by a fixed power of $L_0$.
At successive stages, the right sides involve only the finitely many
already corrected entries and the targets $1/(|S|+1)$. Cramer's rule
therefore leaves a polynomial common denominator after every stage.
Only $\sum_j2^{k-j}=2^k-1$ entries have been altered beyond initial
rounding. Denote the resulting rational array by $r$; it obeys all
constant mixed moments exactly and differs from $r^0$ by
$\operatorname{poly}(m)/L_0$.

\medskip\noindent\emph{Correcting all nonconstant moments by moving the nodes.}
Let $\phi_0,\ldots,\phi_{m-1}$ be the orthonormal shifted Legendre
polynomials, and let $f_a(t)=\int_0^t\phi_{a-1}(u)\,du$ for
$1\le a\le m$. With $r$ fixed, impose
\begin{equation}\label{node}
 \frac1K\sum_q f_a(b_q)\prod_{i\in S}r_{i,q}
 =\int_0^1 f_a(t)t^{|S|}\,dt,
 \quad S\subseteq[k],\ 1\le a\le m.
\end{equation}
Together with the corrected constant equations, these are equivalent
to~\eqref{mom}, because $1,f_1,\ldots,f_m$ span all degree-$m$
polynomials.

At $(b^0,r^0)$ the Jacobian in $b$, multiplied by $K$, evaluates
\[
 \phi_{a-1}(t)\prod_{i\in S}R_i(t),
 \qquad 1\le a\le m, S\subseteq[k].
\]
Its exact empirical Gram is the continuous Gram, since the baseline
quadrature integrates all products through degree $2D$.
The lacunary conditioning theorem gives a polynomial lower eigenvalue.
The number of rows, all evaluation norms, and their first derivatives
are polynomially bounded. Consequently the fixed right inverse of
this baseline Jacobian has polynomial norm from Euclidean residual
norm to maximum node displacement. The rank rounding gives residual
$\operatorname{poly}(m)/L_0$; the change of the Jacobian on a ball of
maximum radius $\rho$ is at most
$\operatorname{poly}(m)(\rho+L_0^{-1}\operatorname{poly}(m))$.

The same fixed-inverse contraction argument as in the reviewed
two-column construction therefore solves~\eqref{node} within
maximum displacement $\operatorname{poly}(m)/L_0$, provided $L_0$
is chosen to be a sufficiently large fixed power of $m+1$.
For clarity, if $A$ is the baseline Jacobian and $B$ a fixed right
inverse, use $b\mapsto b-BF(b,r)$ on
$b^0+\operatorname{range}(B)$. There $BA$ is the identity; the
preceding derivative estimate makes this map a contraction preserving
a ball of twice its initial displacement. No arithmetic condition
is needed on these corrected nodes.

Increase the exponent of $L_0$ so that both rank and node changes
are less than $c_k/(10K)$ and also $o((m+1)^{-T})$.
All ranks and nodes then remain strictly ordered and interior, with
the asserted gaps. Since $|R_i(t)-t|\le\eta$, their mutual closeness
has the required scale. The rational common denominator remains
polynomial by the constant-dimensional correction argument above.
This proves the theorem.
\end{proof}

\paragraph{What this does not encode.}
The products in~\eqref{mom} specify simultaneous threshold events,
not the internal order of the $k$ variables below or above a threshold.
For $k\ge2$, a simultaneous small-dice construction additionally needs
ordered prefix integrals and a positive completion preserving them.
No such implication is asserted here. The theorem isolates a genuine
arithmetic and conditioning subproblem from that remaining order issue.
\end{document}
